\section{Related work}
\label{sec:related}

The algorithms in this section all propagate information about a static parameter of a state-space model, or they are the constructions with which the nested filter is easily confused. They do not share a parameter target. Table~\ref{tab:targets} records the distinction that the rest of the paper uses. The nearest-neighbour row is the target defined in Section~\ref{sec:methods}. The other rows are the algorithms as published.

\begin{table}[htbp]
\centering
\caption{Parameter targets of the neighbouring schemes. The exact posterior is \(\pi_t\). The ensemble Kalman posterior is the posterior induced by the ensemble Kalman likelihood. Citations for the rows are \citet{chopin2012}, \citet{drovandi2022}, \citet{katzfuss2018}, \citet{golightly2026} and \citet{temfack2026}.}
\label{tab:targets}
\small
\begin{tabular}{>{\raggedright\arraybackslash}p{3.15cm}>{\raggedright\arraybackslash}p{4.15cm}>{\raggedright\arraybackslash}p{4.35cm}>{\raggedright\arraybackslash}p{1.7cm}}
\toprule
Method & Weight & Parameter target & Sequential \\
\midrule
SMC2 & Unbiased particle estimate of the increment & Exact posterior \(\pi_t\) & Yes \\
Ensemble MCMC & Ensemble Kalman likelihood & Ensemble Kalman posterior, fixed series & No \\
Particle ensemble Kalman filter & Integrated ensemble Kalman likelihood & Ensemble Kalman posterior & Yes \\
Plain nested filter & Ensemble Kalman increment & \(\hat{\pi}_{\mathrm{enkf}}\) & Yes \\
Rao-Blackwellised weight & Gaussian forecast, true observation density restored & \(\tilde{\pi}_{\mathrm{enkf}}\) & Yes \\
Nearest-neighbour weight & Prediction of the ensemble Kalman increment & \(\pi_{\mathrm{knn}}\) & Yes \\
Ensemble SMC2 & Unbiased Gaussian predictive density & Posterior under that Gaussian model & Yes \\
\bottomrule
\end{tabular}
\end{table}

SMC2 is the scheme whose parameter marginal is exact. \citet[Proposition~1]{chopin2012} write an extended target that includes the particle-filter randomness, and whose \(\theta\)-marginal is \(\pi_t\) for any fixed number of state particles. The argument uses unbiasedness of the product of incremental weights for the true likelihood. Replacing the inner particle filter by an ensemble Kalman filter removes that unbiasedness. The marginal that remains is \(\hat{\pi}_{\mathrm{enkf}}\), not \(\pi_t\). The same proposition therefore gives no distance between \(\hat{\pi}_{\mathrm{enkf}}\) and \(\pi_t\), and no distance between \(\pi_{\mathrm{knn}}\) and \(\hat{\pi}_{\mathrm{enkf}}\). The nearest-neighbour weight is not unbiased for the ensemble Kalman likelihood either.

Ensemble Markov chain Monte Carlo places the ensemble Kalman likelihood in the acceptance ratio of one chain, run on a fixed series \citep{drovandi2022}. One rejuvenation sweep of the nested filter is one iteration of that scheme, but the chain itself is not sequential, and it does not update a posterior as observations arrive. Even when the forecast is exactly Gaussian, the plug-in Gaussian density is a biased estimator of the idealised Gaussian predictive that uses the true forecast mean and covariance. \citet[Section~3.3]{drovandi2022} discuss an unbiased estimator of that Gaussian density, of Ghurye--Olkin type, available when the forecast is Gaussian and the ensemble is large enough. \citet{temfack2026} use the same style of estimator for an epidemic model with a state-dependent observation variance. An estimator of that kind removes the finite-sample bias of the plug-in Gaussian density. It does not remove the error that separates \(\hat{\pi}_{\mathrm{enkf}}\) from \(\pi_t\), because the idealised Gaussian predictive is already an approximation to the true predictive. \citet{temfack2026} are explicit that the Gaussian assumption remains. Their mutation is particle marginal Metropolis--Hastings. They do not use the Rao-Blackwellised weight of Section~\ref{sec:methods}.

A different sequential construction gives the parameter an artificial transition and resamples it with a shrinkage kernel, so that the filter does not collapse onto a single parameter value \citep{liu2001}. The particle ensemble Kalman filter of \citet{katzfuss2018} weights parameter particles by an integrated ensemble Kalman likelihood and, for a constant parameter, needs a kernel of that kind to keep the parameter cloud alive. The nested filter keeps \(\theta\) static inside each rejuvenation and restores diversity by a Metropolis--Hastings move under the likelihood the weight is using. \citet[Proposition~2]{katzfuss2018} gives the variance of the log ensemble Kalman likelihood as being of order \(n/N\) in one linear example, as the observation dimension \(n\) and the ensemble size \(N\) grow, when the forecast members are independent draws from the true forecast. The ensemble inside the nested filter is interacting, so the proposition is an analogy for the cost of controlling that variance. It is not a uniform bound on the log-likelihood, and it is not a bound on \(\pi_{\mathrm{knn}}\).

The weighted ensemble Kalman filter proposes a state ensemble with a Kalman shift and then reweights it \citep{papadakis2010}. The weight contains the observation density and the transition density. \citet{papadakis2010} note that the transition density is usually unavailable when the dynamical time step is finer than the observation interval, which is the regime of the diffusions in Section~\ref{sec:sim}. The Rao-Blackwellised weight in Section~\ref{sec:methods} is built so that the Gaussian forecast appears in the target and in the proposal and therefore cancels. The transition density between observation times is never evaluated. The weighted ensemble Kalman filter is a state filter. It does not supply a posterior for a static parameter, and its weight is not evidence for or against the ordering of \(\hat{\pi}_{\mathrm{enkf}}\) and \(\tilde{\pi}_{\mathrm{enkf}}\).

Delayed acceptance screens a proposal with a cheap approximation and then corrects the screen with the intended target \citep{christen2005}. If the parent kernel is irreducible for that target, the proposal is reversible, and the approximation is positive wherever the proposal is positive, the target remains invariant. \citet[Section~4.2]{golightly2026} use this construction with an inverse-distance neighbour average in place of the ensemble Kalman likelihood at the first stage, and with the ensemble Kalman likelihood at the second stage. The invariant distribution is \(\hat{\pi}_{\mathrm{enkf}}\). \citet[Section~5.1]{sherlock2016} contrast that use with algorithms that take the approximation as the target: draws aimed at the surrogate are not draws from the expensive posterior. \citet{bon2021} put a surrogate inside the mutation of a sequential Monte Carlo sampler and keep the expensive likelihood as the second-stage target. In both cases the second stage is what protects the intended posterior. The nearest-neighbour weight in Section~\ref{sec:methods} is the use those papers set aside. It targets \(\pi_{\mathrm{knn}}\).

The gap left by this literature is specific. \citet{golightly2026} define \(\hat{\pi}_{\mathrm{enkf}}\) and \(\tilde{\pi}_{\mathrm{enkf}}\), compare them empirically, and use the neighbour average as a screen. They do not define the posterior induced by using that average as the weight, they do not separate that error from the ensemble Kalman error, and they do not record when the Rao-Blackwellised weight collapses to the plain weight. The unbiased Gaussian correction in \citet{drovandi2022} and \citet{temfack2026} addresses a third error, the plug-in bias of a Gaussian density, and leaves both of the errors in Figure~\ref{fig:targets} untouched. Section~\ref{sec:methods} treats the two errors that the nested filter can make. Section~\ref{sec:sim} compares the targets without adding the errors together.
